% ===================================================================== Chapter 3
\chapter{Theoretical Foundation}\label{ch:theory}
This chapter collects the relations used by the analytical baseline and for interpreting the simulations. They are standard published
theory; what is specific to this study is the choice of relations and the parameters fed into them. SI units are used throughout.

\section{Internal Flow}
Fluent solves the steady Reynolds-averaged conservation equations for mass, momentum and energy in the air, closed with the $k$--$\omega$
SST model \cite{menter1994}, and the steady heat-conduction equation $\nabla\cdot(k_s\nabla T) = 0$ in the solid. The two regions are
coupled by continuity of temperature and heat flux at the interface. The air is treated as an incompressible ideal gas,
\begin{equation}
\rho = \frac{p_{op}}{R\,T}, \qquad R = 287.058~\text{J/(kg\,K)},\; p_{op} = 101{,}325~\text{Pa},
\label{eq:idealgas}
\end{equation}
so that density follows temperature but not the local pressure; this is admissible because the pressure drop is 0.43\,\% of $p_{op}$.

\section{Reynolds Number}
With mass flux $G = \dot m/A_c$, which stays constant along the duct although the density falls,
\begin{equation}
Re = \frac{\rho V D_h}{\mu} = \frac{G D_h}{\mu}, \qquad D_h = \frac{4A_c}{P_w} = D_i ,
\label{eq:re}
\end{equation}
where $V$ is the mean velocity [m/s], $\mu$ the dynamic viscosity [Pa\,s] and $P_w$ the wetted perimeter [m]. Flow in a circular duct is
turbulent for $Re > 4000$.

\section{Pressure Loss}
The static pressure drop of the heated duct consists of wall friction, thermal acceleration of the expanding gas and the extra loss of the
developing entry region:
\begin{align}
\Delta p_{fr} &= f\,\frac{L}{D_h}\,\frac{\rho V^2}{2}, \label{eq:darcy}\\
f &= \left(0.790\ln Re - 1.64\right)^{-2} \quad (3\times10^3 < Re < 5\times10^6), \label{eq:petukhov}\\
\Delta p_{acc} &= G^2\left(\frac{1}{\rho_{out}} - \frac{1}{\rho_{in}}\right), \qquad \Delta p_{ent} = K_\infty\,\frac{\rho V^2}{2}, \label{eq:acc}
\end{align}
with the Darcy friction factor $f$ from Petukhov \cite{petukhov1970} and the turbulent entry increment $K_\infty \approx 0.09$
\cite{kayslondon}. Entrance and exit losses of a real installation are not part of the CFD domain and are never compared with Fluent.

\section{Turbulent Heat Transfer}
The heat-transfer coefficient $h$ [W/(m$^2$\,K)] is defined by Newton's law of cooling,
\begin{equation}
q''_i = h\,(T_{w,i} - T_b),
\label{eq:newton}
\end{equation}
where $q''_i$ is the heat flux at the bore, $T_{w,i}$ the inner-wall temperature and $T_b$ the local bulk temperature. Because the same
heat passes through a smaller area, the bore flux is $q''_i = q''_o D_o/D_i$ = 16,000~W/m$^2$ for the baseline.

\section{Nusselt Number}
The Nusselt number and the two correlations used are
\begin{align}
Nu &= \frac{h D_h}{k}, \label{eq:nu}\\
Nu_{Gn} &= \frac{(f/8)(Re-1000)Pr}{1 + 12.7\sqrt{f/8}\,\left(Pr^{2/3}-1\right)} \quad (3\times10^3 < Re < 5\times10^6,\; 0.5 < Pr < 2000), \label{eq:gnielinski}\\
Nu_{DB} &= 0.023\,Re^{0.8}Pr^{0.4}, \label{eq:db}
\end{align}
from Gnielinski \cite{gnielinski1976} and Dittus and Boelter \cite{dittus1930}. Both are constant-property correlations; for a gas heated
with a large wall-to-bulk temperature difference the analytical model applies the property-ratio correction \cite{kayscrawford}
\begin{equation}
Nu = Nu_{cp}\left(\frac{T_b}{T_w}\right)^{n}, \qquad n \approx 0.5 .
\label{eq:propcorr}
\end{equation}

\section{Convective Heat Transfer}
Equations~\eqref{eq:newton}--\eqref{eq:propcorr} give the film temperature difference $T_{w,i} - T_b = q''_i/h$. In the analytical model this
difference is about 206~K at the exit, far larger than the 7~K through the wall. In the CFD, $h$ follows from the resolved near-wall
temperature gradient; the ratio of the CFD value to Equation~\eqref{eq:propcorr} is itself a result (Chapter~\ref{ch:cfdres}).

\section{Conduction Through Solid Wall}
For steady radial conduction without generation, with heat per unit length $q'$ [W/m],
\begin{align}
T(r) &= T(r_i) + \frac{q'}{2\pi k_s}\ln\frac{r}{r_i}, \qquad \Delta T_{wall} = \frac{q'}{2\pi k_s}\ln\frac{r_o}{r_i}, \label{eq:fourier}\\
R_{cond} &= \frac{\ln(r_o/r_i)}{2\pi k_s L}, \qquad R_{conv} = \frac{1}{h A_i}. \label{eq:resist}
\end{align}
The temperature rises logarithmically towards the heated outer surface. The resistance ratio $R_{cond}/R_{conv}$ decides whether the metal
or the air film controls the wall temperature.

\section{Energy Balance}
The imposed heat is carried away by the air:
\begin{equation}
Q = q''_o\,A_o = \dot m\,c_p\,(T_{out} - T_{in}), \qquad \frac{dT_b}{dz} = \frac{q''_i\,P_w}{\dot m\,c_p} .
\label{eq:energy}
\end{equation}
At constant flux the bulk temperature rises linearly along the duct. The outlet bulk temperature depends only on $Q$, $\dot m$ and $c_p$,
not on $h$, which makes Equation~\eqref{eq:energy} the strongest single check of a CFD solution.

\section{Thermal Expansion}
The free thermal strain and the resulting free axial growth are
\begin{equation}
\varepsilon_{th} = \alpha_s(T)\,\left(T - T_{ref}\right), \qquad \Delta L = \int_0^L \bar\varepsilon_{th}(z)\,dz ,
\label{eq:thermalstrain}
\end{equation}
where $\alpha_s$ is the secant (mean) expansion coefficient referred to $T_{ref}$ = 300~K and $\bar\varepsilon_{th}$ the section-mean
strain. Thermoelastic Hooke's law adds $\varepsilon_{th}$ to the mechanical strain:
$\varepsilon_{ij} = \left[(1+\nu)\sigma_{ij} - \nu\sigma_{kk}\delta_{ij}\right]/E + \varepsilon_{th}\,\delta_{ij}$.

\section{Thermal Stress}
Stress appears only where thermal strain is prevented, by external restraint or by neighbouring material at a different temperature. For a
fully restrained member with a section-mean temperature rise $\Delta\bar T$,
\begin{equation}
\bar\sigma_z = -E\,\alpha_s\,\Delta\bar T, \qquad N = \int_A E\,\alpha_s\,\Delta T\,dA .
\label{eq:restrained}
\end{equation}
For a long cylinder with free ends and a logarithmic radial temperature distribution, the hoop stress at the bore is
\cite{timoshenko}
\begin{equation}
\sigma_\theta(a) = \frac{\alpha E\,T_a}{2(1-\nu)\ln(b/a)}\left[1 - \frac{2b^2}{b^2-a^2}\ln\frac{b}{a}\right], \qquad T_a = T_i - T_o ,
\label{eq:timoshenko}
\end{equation}
with $a = r_i$ and $b = r_o$. Stresses are assessed with the von Mises equivalent stress
\begin{equation}
\sigma_{vM} = \sqrt{\tfrac12\left[(\sigma_r-\sigma_\theta)^2 + (\sigma_\theta-\sigma_z)^2 + (\sigma_z-\sigma_r)^2\right]}
\label{eq:vm}
\end{equation}
against the yield strength at the local temperature, which defines the yield utilisation $\sigma_{vM}/S_y(T)$ and the first-yield load
factor $S_y(T)/\sigma_{vM}$.

\section{Buckling Fundamentals}
A straight column under axial compression $N$ has the elastic critical load
\begin{equation}
P_E = \frac{\pi^2 E I}{(K L)^2}, \qquad r = \sqrt{I/A},
\label{eq:euler}
\end{equation}
where $K$ is the effective-length factor set by the end conditions: 1.0 for a guided column whose ends can translate but not rotate,
0.699 for fixed--pinned and 0.5 for fixed--fixed. For a stocky section the shear deformation reduces the load,
$P = P_E/\left(1 + P_E/(k G A)\right)$ with Cowper's shear coefficient $k$ \cite{cowper1966}. When the slenderness $KL/r$ is below
$C_c = \sqrt{2\pi^2E/S_y}$ the column is intermediate and the conventional Johnson parabola
$\sigma_{cr} = S_y - \left(S_y^2/4\pi^2E\right)(KL/r)^2$ gives an inelastic estimate. Because the axial force here is thermal and
displacement-controlled, the buckling load factor is expressed as $P_{cr}/N$.

\section{Linear Eigenvalue Buckling}
ANSYS computes the load factor $\lambda$ that multiplies the pre-stress state of LC2 (the full thermal load) from the eigenvalue problem
\begin{equation}
\left([K] + \lambda\,[K_\sigma]\right)\{\phi\} = \{0\}, \qquad P_{cr} = \lambda_1\,N ,
\label{eq:eigen}
\end{equation}
where $[K]$ is the elastic stiffness matrix, $[K_\sigma]$ the stress-stiffness matrix of the LC2 solution and $\{\phi\}$ the mode shape. The
first eigenvalue $\lambda_1$ is the bifurcation load factor of a perfectly straight, linear-elastic structure with idealised supports. It
contains no imperfection, no plasticity and no large-deflection effect. It is therefore a linear stability indicator for an idealised
model, \textbf{not} a guaranteed real-world collapse load and not a factor of safety; an imperfect real duct deflects laterally before
$\lambda_1$ is reached and its capacity is lower by an amount that only a nonlinear analysis can quantify.
